
\subsubsection{6.1 Interpretation}

Results validate that architectural differences decompose into three learnable signals, each captured through specialized components. All three contributed independently (ablation: +0.04, +0.05, +0.04) without degrading one another.

The hierarchy observed---contrastive projection (0.05) > operation encoding (0.04) $\approx$ graph topology (0.04)---suggests prioritization for future work. Contrastive projection creates the metric space foundation enabling cross-architecture comparison, while operation encoders and GNN provide structural and topological refinements.

The high GNN residual weight ($\alpha$ = 0.42) indicates computational graph topology encodes property-predictive information beyond operation types. ResNet's skip connections create gradient highways affecting optimization stability \citep{He2016}, while ViT's dense attention creates all-to-all communication within blocks affecting information integration. These topological differences appear learnable and informative.

\subsubsection{6.2 Limitations}

\textbf{Proof-of-Concept Scale:} Experiments used 100 models (50 ResNet-50, 50 ViT-Base) with 10 training epochs. Full-scale validation with 400 models across 4 architectures (ResNet, ViT, MobileNet, EfficientNet) and 100 epochs is required to test generalization. Claims are restricted to ResNet-50 and ViT-Base until broader architectural diversity is tested.

\textbf{Synthetic Accuracy Labels:} H-M-Integrated used synthetic ImageNet accuracy labels sampled from normal distributions (mean/std from published model card statistics) rather than real top-1 accuracy. This enabled mechanism validation---all diagnostic falsifiers passed---but property prediction capability is not demonstrated. The measured $\rho$ = 0.67 reflects learned patterns in synthetic distributions, not genuine accuracy prediction from weights. H-E1 binary classifier independently validated that operation-specific signals exist in real pretrained weights (91.7\% accuracy distinguishing ResNet from ViT with mock data), providing partial evidence of real-weight signal existence.

\textbf{Vision Domain Only:} Results hold for ImageNet-1K image classification models across convolutional and transformer architectures. Extension to NLP or speech requires adapting operation encoders. Contrastive projection and GNN residual are domain-agnostic but remain empirically untested outside vision.

\textbf{Single Transfer Direction:} Proof-of-concept focused on $ResNet\rightarrow ViT$ transfer. Reverse direction ($ViT\rightarrow ResNet)$ and other pairs remain untested. Within-architecture transfer ($ResNet\rightarrow ResNet, ViT\rightarrow ViT)$ was not measured. Full transfer matrix validation is required.

\subsubsection{6.3 Comparison to Prior Work}

\textbf{SNE ($\rho$ = 0.54):} Achieved cross-architecture capability through architecture-invariance, discarding operation-specific structure via uniform tokenization. CAPE's +0.13 improvement demonstrates that architecture-conditioning---preserving operation specificity, task alignment, and graph topology---outperforms invariance.

\textbf{SANE:} Demonstrated operation-specific encoding improves within-architecture transfer (+2.2\% accuracy for ResNet-$18\rightarrow ResNet-50)$ but required shared token sizes. CAPE extends operation-specific encoding to cross-architecture transfer via contrastive alignment, removing the same-family constraint.

\textbf{UNF:} Provided theoretical foundations for permutation-equivariant weight-space learning. CAPE applies UNF's equivariance to attention encoders within a practical cross-architecture system, providing empirical validation at proof-of-concept scale.

\subsubsection{6.4 Future Work}

Full-scale validation with 400 models across 4 architectures and real ImageNet accuracy labels is required to demonstrate predictive capability beyond mechanism validation. Additional transfer directions ($ViT\rightarrow ResNet, ResNet\rightarrow MobileNet)$ should be tested. Extension to multi-task model zoos (ImageNet + CIFAR-10 + MS-COCO) would test contrastive alignment generalization. Scaling to larger models (>100M parameters) likely requires no architectural changes, only increased computational resources.
